\section{AHSME 1994}
        \begin{enumerate}
        	\item (\textit{AHSME 1994}) $4^4 \cdot 9^4 \cdot 4^9 \cdot 9^9=$


 $\textbf{(A)}\ 13^{13} \qquad\textbf{(B)}\ 13^{36} \qquad\textbf{(C)}\ 36^{13} \qquad\textbf{(D)}\ 36^{36} \qquad\textbf{(E)}\ 1296^{26}$


 

	\item (\textit{AHSME 1994}) A large rectangle is partitioned into four rectangles by two segments parallel to its sides. The areas of three of the resulting rectangles are shown. What is the area of the fourth rectangle?

                    \begin{figure}[H]
                        \centering
                        \includegraphics[width=0.5\textwidth]{Problem/AHSME_I_1994_2_Problem_Figure_1.png}
                    \end{figure}
                    


 $\textbf{(A)}\ 10 \qquad\textbf{(B)}\ 15 \qquad\textbf{(C)}\ 20 \qquad\textbf{(D)}\ 21 \qquad\textbf{(E)}\ 25$


 

	\item (\textit{AHSME 1994}) How many of the following are equal to $x^x+x^x$ for all $x>0$?


 $\textbf{I:}\ 2x^x \qquad\textbf{II:}\ x^{2x} \qquad\textbf{III:}\ (2x)^x \qquad\textbf{IV:}\ (2x)^{2x}$


 $\textbf{(A)}\ 0 \qquad\textbf{(B)}\ 1 \qquad\textbf{(C)}\ 2 \qquad\textbf{(D)}\ 3 \qquad\textbf{(E)}\ 4$


 

	\item (\textit{AHSME 1994}) In the $xy$-plane, the segment with endpoints $(-5,0)$ and $(25,0)$ is the diameter of a circle. If the point $(x,15)$ is on the circle, then $x=$


 $\textbf{(A)}\ 10 \qquad\textbf{(B)}\ 12.5 \qquad\textbf{(C)}\ 15\qquad\textbf{(D)}\ 17.5 \qquad\textbf{(E)}\ 20$


 

	\item (\textit{AHSME 1994}) Pat intended to multiply a number by $6$ but instead divided by $6$. Pat then meant to add $14$ but instead subtracted $14$. After these mistakes, the result was $16$. If the correct operations had been used, the value produced would have been


 $\textbf{(A)}\ \text{less than 400} \qquad\textbf{(B)}\ \text{between 400 and 600} \qquad\textbf{(C)}\ \text{between 600 and 800} \\ \textbf{(D)}\ \text{between 800 and 1000} \qquad\textbf{(E)}\ \text{greater than 1000}$


 

	\item (\textit{AHSME 1994}) In the sequence

 \[..., a, b, c, d, 0, 1, 1, 2, 3, 5, 8,...\]
each term is the sum of the two terms to its left. Find $a$.


 $\textbf{(A)}\ -3 \qquad\textbf{(B)}\ -1 \qquad\textbf{(C)}\ 0 \qquad\textbf{(D)}\ 1 \qquad\textbf{(E)}\ 3$


 

	\item (\textit{AHSME 1994}) Squares $ABCD$ and $EFGH$ are congruent, $AB=10$, and $G$ is the center of square $ABCD$. The area of the region in the plane covered by these squares is

                    \begin{figure}[H]
                        \centering
                        \includegraphics[width=0.5\textwidth]{Problem/AHSME_I_1994_7_Problem_Figure_1.png}
                    \end{figure}
                    
$\textbf{(A)}\ 75 \qquad\textbf{(B)}\ 100 \qquad\textbf{(C)}\ 125 \qquad\textbf{(D)}\ 150 \qquad\textbf{(E)}\ 175$


 

	\item (\textit{AHSME 1994}) In the polygon shown, each side is perpendicular to its adjacent sides, and all 28 of the sides are congruent. The perimeter of the polygon is $56$. The area of the region bounded by the polygon is

                    \begin{figure}[H]
                        \centering
                        \includegraphics[width=0.5\textwidth]{Problem/AHSME_I_1994_8_Problem_Figure_1.png}
                    \end{figure}
                    
$\textbf{(A)}\ 84 \qquad\textbf{(B)}\ 96 \qquad\textbf{(C)}\ 100 \qquad\textbf{(D)}\ 112 \qquad\textbf{(E)}\ 196$


 

	\item (\textit{AHSME 1994}) If $\angle A$ is four times $\angle B$, and the complement of $\angle B$ is four times the complement of $\angle A$, then $\angle B=$


 $\textbf{(A)}\ 10^{\circ} \qquad\textbf{(B)}\ 12^{\circ} \qquad\textbf{(C)}\ 15^{\circ} \qquad\textbf{(D)}\ 18^{\circ} \qquad\textbf{(E)}\ 22.5^{\circ}$


 

	\item (\textit{AHSME 1994}) For distinct real numbers $x$ and $y$, let $M(x,y)$ be the larger of $x$ and $y$ and let $m(x,y)$ be the smaller of $x$ and $y$. If $a<b<c<d<e$, then

 \[M(M(a,m(b,c)),m(d,m(a,e)))=\]
$\textbf{(A)}\ a \qquad\textbf{(B)}\ b \qquad\textbf{(C)}\ c \qquad\textbf{(D)}\ d \qquad\textbf{(E)}\ e$


 

	\item (\textit{AHSME 1994}) Three cubes of volume $1, 8$ and $27$ are glued together at their faces. The smallest possible surface area of the resulting configuration is


 $\textbf{(A)}\ 36 \qquad\textbf{(B)}\ 56 \qquad\textbf{(C)}\ 70 \qquad\textbf{(D)}\ 72 \qquad\textbf{(E)}\ 74$


 

	\item (\textit{AHSME 1994}) If $i^2=-1$, then $(i-i^{-1})^{-1}=$


 $\textbf{(A)}\ 0 \qquad\textbf{(B)}\ -2i \qquad\textbf{(C)}\ 2i \qquad\textbf{(D)}\ -\frac{i}{2} \qquad\textbf{(E)}\ \frac{i}{2}$


 

	\item (\textit{AHSME 1994}) In triangle $ABC$, $AB=AC$. If there is a point $P$ strictly between $A$ and $B$ such that $AP=PC=CB$, then $\angle A =$

                    \begin{figure}[H]
                        \centering
                        \includegraphics[width=0.5\textwidth]{Problem/AHSME_I_1994_13_Problem_Figure_1.png}
                    \end{figure}
                    
$\textbf{(A)}\ 30^{\circ} \qquad\textbf{(B)}\ 36^{\circ} \qquad\textbf{(C)}\ 48^{\circ} \qquad\textbf{(D)}\ 60^{\circ} \qquad\textbf{(E)}\ 72^{\circ}$


 

	\item (\textit{AHSME 1994}) Find the sum of the arithmetic series

 \[20+20\frac{1}{5}+20\frac{2}{5}+\cdots+40\]
$\textbf{(A)}\ 3000 \qquad\textbf{(B)}\ 3030 \qquad\textbf{(C)}\ 3150 \qquad\textbf{(D)}\ 4100 \qquad\textbf{(E)}\ 6000$


 

	\item (\textit{AHSME 1994}) For how many $n$ in $\{1, 2, 3, ..., 100 \}$ is the tens digit of $n^2$ odd?


 $\textbf{(A)}\ 10 \qquad\textbf{(B)}\ 20 \qquad\textbf{(C)}\ 30 \qquad\textbf{(D)}\ 40 \qquad\textbf{(E)}\ 50$


 

	\item (\textit{AHSME 1994}) Some marbles in a bag are red and the rest are blue. If one red marble is removed, then one-seventh of the remaining marbles are red. If two blue marbles are removed instead of one red, then one-fifth of the remaining marbles are red. How many marbles were in the bag originally?


 $\textbf{(A)}\ 8 \qquad\textbf{(B)}\ 22 \qquad\textbf{(C)}\ 36 \qquad\textbf{(D)}\ 57 \qquad\textbf{(E)}\ 71$


 

	\item (\textit{AHSME 1994}) An $8$ by $2\sqrt{2}$ rectangle has the same center as a circle of radius $2$. The area of the region common to both the rectangle and the circle is


 $\textbf{(A)}\ 2\pi \qquad\textbf{(B)}\ 2\pi+2 \qquad\textbf{(C)}\ 4\pi-4 \qquad\textbf{(D)}\ 2\pi+4 \qquad\textbf{(E)}\ 4\pi-2$


 

	\item (\textit{AHSME 1994}) Triangle $ABC$ is inscribed in a circle, and $\angle B = \angle C = 4\angle A$. If $B$ and $C$ are adjacent vertices of a regular polygon of $n$ sides inscribed in this circle, then $n=$

                    \begin{figure}[H]
                        \centering
                        \includegraphics[width=0.5\textwidth]{Problem/AHSME_I_1994_18_Problem_Figure_1.png}
                    \end{figure}
                    
$\textbf{(A)}\ 5 \qquad\textbf{(B)}\ 7 \qquad\textbf{(C)}\ 9 \qquad\textbf{(D)}\ 15 \qquad\textbf{(E)}\ 18$


 

	\item (\textit{AHSME 1994}) Label one disk "$1$", two disks "$2$", three disks "$3$"$, ...,$ fifty disks "$50$". Put these $1+2+3+ \cdots+50=1275$ labeled disks in a box. Disks are then drawn from the box at random without replacement. The minimum number of disks that must be drawn to guarantee drawing at least ten disks with the same label is


 $\textbf{(A)}\ 10 \qquad\textbf{(B)}\ 51 \qquad\textbf{(C)}\ 415 \qquad\textbf{(D)}\ 451 \qquad\textbf{(E)}\ 501$


 

	\item (\textit{AHSME 1994}) Suppose $x,y,z$ is a geometric sequence with common ratio $r$ and $x \neq y$. If $x, 2y, 3z$ is an arithmetic sequence, then $r$ is


 $\textbf{(A)}\ \frac{1}{4} \qquad\textbf{(B)}\ \frac{1}{3} \qquad\textbf{(C)}\ \frac{1}{2} \qquad\textbf{(D)}\ 2 \qquad\textbf{(E)}\ 4$


 

	\item (\textit{AHSME 1994}) Find the number of counter examples to the statement:

 \[``\text{If  N is an odd positive integer the sum of whose digits is 4 and none of whose digits is 0, then N is prime}."\]
$\textbf{(A)}\ 0 \qquad\textbf{(B)}\ 1 \qquad\textbf{(C)}\ 2 \qquad\textbf{(D)}\ 3 \qquad\textbf{(E)}\ 4$


 

	\item (\textit{AHSME 1994}) Nine chairs in a row are to be occupied by six students and Professors Alpha, Beta and Gamma. These three professors arrive before the six students and decide to choose their chairs so that each professor will be between two students. In how many ways can Professors Alpha, Beta and Gamma choose their chairs?


 $\textbf{(A)}\ 12 \qquad\textbf{(B)}\ 36 \qquad\textbf{(C)}\ 60 \qquad\textbf{(D)}\ 84 \qquad\textbf{(E)}\ 630$


 

	\item (\textit{AHSME 1994}) In the $xy$-plane, consider the L-shaped region bounded by horizontal and vertical segments with vertices at $(0,0), (0,3), (3,3), (3,1), (5,1)$ and $(5,0)$. The slope of the line through the origin that divides the area of this region exactly in half is

                    \begin{figure}[H]
                        \centering
                        \includegraphics[width=0.5\textwidth]{Problem/AHSME_I_1994_23_Problem_Figure_1.png}
                    \end{figure}
                    
$\textbf{(A)}\ \frac{2}{7} \qquad\textbf{(B)}\ \frac{1}{3} \qquad\textbf{(C)}\ \frac{2}{3} \qquad\textbf{(D)}\ \frac{3}{4} \qquad\textbf{(E)}\ \frac{7}{9}$


 

	\item (\textit{AHSME 1994}) A sample consisting of five observations has an arithmetic mean of $10$ and a median of $12$. The smallest value that the range (largest observation minus smallest) can assume for such a sample is


 $\textbf{(A)}\ 2 \qquad\textbf{(B)}\ 3 \qquad\textbf{(C)}\ 5 \qquad\textbf{(D)}\ 7 \qquad\textbf{(E)}\ 10$


 

	\item (\textit{AHSME 1994}) If $x$ and $y$ are non-zero real numbers such that

 \[|x|+y=3 \qquad \text{and} \qquad |x|y+x^3=0,\]
then the integer nearest to $x-y$ is


 $\textbf{(A)}\ -3 \qquad\textbf{(B)}\ -1 \qquad\textbf{(C)}\ 2 \qquad\textbf{(D)}\ 3 \qquad\textbf{(E)}\ 5$


 

	\item (\textit{AHSME 1994}) A regular polygon of $m$ sides is exactly enclosed (no overlaps, no gaps) by $m$ regular polygons of $n$ sides each. (Shown here for $m=4, n=8$.) If $m=10$, what is the value of $n$?

                    \begin{figure}[H]
                        \centering
                        \includegraphics[width=0.5\textwidth]{Problem/AHSME_I_1994_26_Problem_Figure_1.png}
                    \end{figure}
                    
$\textbf{(A)}\ 5 \qquad\textbf{(B)}\ 6 \qquad\textbf{(C)}\ 14 \qquad\textbf{(D)}\ 20 \qquad\textbf{(E)}\ 26$


 

	\item (\textit{AHSME 1994}) A bag of popping corn contains $\frac{2}{3}$ white kernels and $\frac{1}{3}$ yellow kernels. Only $\frac{1}{2}$ of the white kernels will pop, whereas $\frac{2}{3}$ of the yellow ones will pop. A kernel is selected at random from the bag, and pops when placed in the popper. What is the probability that the kernel selected was white?


 $\textbf{(A)}\ \frac{1}{2} \qquad\textbf{(B)}\ \frac{5}{9} \qquad\textbf{(C)}\ \frac{4}{7} \qquad\textbf{(D)}\ \frac{3}{5} \qquad\textbf{(E)}\ \frac{2}{3}$


 

	\item (\textit{AHSME 1994}) In the $xy$-plane, how many lines whose $x$-intercept is a positive prime number and whose $y$-intercept is a positive integer pass through the point $(4,3)$?


 $\textbf{(A)}\ 0 \qquad\textbf{(B)}\ 1 \qquad\textbf{(C)}\ 2 \qquad\textbf{(D)}\ 3 \qquad\textbf{(E)}\ 4$


 

	\item (\textit{AHSME 1994}) Points $A, B$ and $C$ on a circle of radius $r$ are situated so that $AB=AC, AB>r$, and the length of minor arc $BC$ is $r$. If angles are measured in radians, then $AB/BC=$

                    \begin{figure}[H]
                        \centering
                        \includegraphics[width=0.5\textwidth]{Problem/AHSME_I_1994_29_Problem_Figure_1.png}
                    \end{figure}
                    
$\textbf{(A)}\ \frac{1}{2}\csc{\frac{1}{4}} \qquad\textbf{(B)}\ 2\cos{\frac{1}{2}} \qquad\textbf{(C)}\ 4\sin{\frac{1}{2}} \qquad\textbf{(D)}\ \csc{\frac{1}{2}} \qquad\textbf{(E)}\ 2\sec{\frac{1}{2}}$


 

	\item (\textit{AHSME 1994}) When $n$ standard 6-sided dice are rolled, the probability of obtaining a sum of 1994 is greater than zero and is the same as the probability of obtaining a sum of $S$. The smallest possible value of $S$ is


 $\textbf{(A)}\ 333 \qquad\textbf{(B)}\ 335 \qquad\textbf{(C)}\ 337 \qquad\textbf{(D)}\ 339 \qquad\textbf{(E)}\ 341$


 

\end{enumerate}